\chapter{Conclusão e Trabalhos Futuros}
\label{cap:conclusao}

\section{Síntese do Trabalho}
\label{sec:sintese}

Esta dissertação investigou o problema da janela de contexto em Modelos de Linguagem de Grande Escala a partir de uma perspectiva não invasiva, isto é, aplicável a modelos acessados via API, sem retreinamento nem modificação arquitetural. O trabalho partiu da constatação de que janelas de contexto nominalmente grandes frequentemente não são utilizadas de forma eficaz pelos modelos, em razão da complexidade quadrática da atenção e de fenômenos como o viés \textit{lost in the middle}, o \textit{context rot} e o colapso de atenção.

O Capítulo~\ref{cap:fundamentos} revisou os fundamentos técnicos dos LLMs e organizou as estratégias não invasivas de extensão de contexto em uma taxonomia, composta por compressão de \textit{prompt}, segmentação de contexto e janela deslizante, métodos baseados em \textit{MapReduce} e alocação de memória externa. A revisão identificou como lacuna recorrente o tratamento uniforme do \textit{prompt}, em que um único método e uma única taxa de compressão são aplicados a toda a entrada, desconsiderando sua estrutura interna heterogênea.

O Capítulo~\ref{cap:metodologia} apresentou a validação empírica das estratégias não invasivas em cinco modelos de arquiteturas distintas, que reúnem redes densas, redes de mistura de especialistas e modelos ajustados para raciocínio, sobre uma suíte de \textit{benchmarks} de contexto longo. A estratégia \textit{Sliding Window} apresentou o melhor equilíbrio recorrente entre acurácia e latência na etapa inicial, o que sustentou sua adoção como estratégia base. A triagem adicional dos compressores abrangeu os cinco modelos e os dois orçamentos, totalizando 880 execuções e 879 respostas válidas. CPC-MiniLM obteve as maiores médias entre os quatro compressores em oito das dez combinações, enquanto LLMLingua-2 liderou nas duas combinações do Qwen 3 30B-A3B. Selective Context, por sua vez, apresentou a maior fidelidade semântica segundo a métrica adotada. Esses perfis complementares sustentaram a incorporação dos três ao conjunto inicial de opções do MCKP. Em seguida, o mesmo capítulo formalizou o \textit{framework} proposto, que descreve a compressão seletiva de contexto como um problema da mochila de múltipla escolha, com função de importância, modelo de qualidade, restrições de compatibilidade e solução por programação dinâmica.

O Capítulo~\ref{cap:resultados} discutiu os resultados empíricos, analisou a relação entre qualidade, compressão e latência e delimitou as implicações da triagem para o conjunto de opções do \textit{framework}. O mesmo capítulo apresentou a avaliação do resolvedor MCKP no Llama 3.1 8B em duas janelas de contexto. Quando o contexto excede o orçamento, a compressão supera o envio truncado, e sob a janela de 2048 \textit{tokens} a alocação seletiva alcança 0.416 nesses casos, contra 0.347 da compressão uniforme. Em uma amostra ampliada de vinte casos por \textit{benchmark}, a vantagem do MCKP sobre o truncamento é estatisticamente significativa, com Wilcoxon $p = 0.015$, enquanto a vantagem sobre a compressão uniforme permanece positiva, porém marginal, com $p = 0.052$. Essa vantagem aparece somente na janela de 2048 \textit{tokens}, e não na de 8192, o que delimita seu alcance ao regime de orçamento restritivo.

\section{Contribuições}
\label{sec:contribuicoes-conclusao}

O trabalho ofereceu quatro contribuições principais. A primeira é uma revisão sistemática de técnicas não invasivas de extensão de contexto, organizada em uma taxonomia por lógica operacional. A segunda é uma avaliação empírica comparativa conduzida em duas fases, uma preliminar sobre modelos servidos por APIs externas e outra local, que compara nove técnicas de tratamento de contexto em cinco modelos de arquiteturas distintas e dois orçamentos, reportando desempenho, latência e taxa retida, com a fidelidade semântica medida na etapa de ampliação dedicada aos quatro compressores. A terceira é a formalização de um \textit{framework} de compressão seletiva via problema da mochila, que atua como uma camada de decisão acima de compressores existentes. A quarta é uma avaliação do \textit{framework} sob dois orçamentos de contexto, que indica que a alocação seletiva por partição supera o truncamento de forma estatisticamente significativa e a compressão uniforme por margem positiva ainda não significativa quando o contexto excede a janela do modelo, acompanhada de um plano experimental reprodutível, orientado por métricas de preservação de informação, para as configurações ainda não executadas.

\section{Limitações}
\label{sec:limitacoes}

O trabalho apresenta limitações que devem ser consideradas na interpretação dos resultados. A triagem dos quatro compressores utiliza 22 casos por combinação, sem repetições independentes, e uma chamada do CPC-MiniLM no DeepSeek-R1 32B sob 8192 \textit{tokens} terminou por \textit{timeout}; essa combinação possui 21 respostas válidas. A fidelidade semântica adicionada à triagem mede cobertura por \textit{embeddings}, mas não captura integralmente o risco de alucinação por perda de conteúdo. No caso de CPC-MiniLM, essa métrica usa o mesmo codificador empregado pela aproximação, não constituindo uma avaliação inteiramente independente. Além disso, a verificação empírica do \textit{framework} restringe-se ao Llama 3.1 8B e a duas janelas de contexto. Em uma amostra de vinte casos por \textit{benchmark}, a hipótese H1 recebe suporte parcial, uma vez que a vantagem do MCKP é significativa frente ao truncamento, com $p = 0.015$, e apenas marginal frente à compressão uniforme, com $p = 0.052$. As hipóteses H2 e H3, que dependem de estudos de ablação da saliência posicional e da penalização de transição, permanecem a ser testadas. A avaliação também identificou um desalinhamento entre a fidelidade medida pelo codificador MiniLM e o desempenho final, que penaliza a seleção de opções e ajuda a explicar por que a vantagem sobre a compressão uniforme ainda não é conclusiva.

\section{Trabalhos Futuros}
\label{sec:trabalhos-futuros}

A partir das limitações identificadas, delineiam-se as seguintes direções de trabalho futuro. A primeira é ampliar a avaliação do resolvedor iniciada no Capítulo~\ref{cap:resultados}, estendendo-a aos demais modelos e orçamentos e aumentando o número de casos, de modo a confirmar a vantagem sobre a compressão uniforme, que nesta dissertação permaneceu marginal, e executando os estudos de ablação da saliência posicional e da penalização de transição associados às hipóteses H2 e H3. A segunda é alinhar a função de valor à qualidade final, substituindo a fidelidade por cobertura de \textit{embeddings} por uma medida mais correlacionada com o desempenho \textit{downstream}, e incorporar métricas adicionais de preservação de informação, preferencialmente independentes dos codificadores usados pelos compressores. A terceira é avaliar CPC com os \textit{checkpoints} disponibilizados pelo trabalho original, separando o custo dessa configuração dos resultados da aproximação CPC-MiniLM. A quarta é estender o \textit{framework} a modelos multimodais, em que partições podem conter texto, imagens ou outros tipos de conteúdo, com compressores específicos por modalidade. A quinta é otimizar a latência do próprio \textit{framework}, uma vez que as etapas de estimativa de importância e de geração de opções introduzem custo computacional que precisa ser compatível com aplicações interativas.

De forma mais ampla, o trabalho sugere que a compressão de contexto em LLMs pode ser tratada como um problema de alocação de recursos, e não como uma redução uniforme de \textit{tokens}. Essa formulação abre espaço para políticas de compressão adaptativas, que respondem à estrutura interna do \textit{prompt} e ao orçamento disponível, e que podem ser aplicadas de forma não invasiva sobre modelos já treinados.
